\chapter{Starter Listings}

These are the starter listings currently stored on \diskname. Each one calls
the common framework and leaves a clear place for the full game code. The short
introductions are here to answer the most important question before anyone
starts typing: why is this program worth bringing to life?

\programlisting{Tiny 1}{tiny-1.bas}{ball}{This is the quickest way to see NovaBASIC make something feel alive. A shiny hardware sprite falls, speeds up under gravity, rebounds inside the visible screen border, and waits for the video frame so the motion feels smooth instead of jumpy. In just a few lines, you get physics, sprites, sound, and timing working together.}
\programlisting{Tiny 2}{tiny-2.bas}{color}{Color is one of the easiest ways to make a small program feel bigger than it is. This tiny demo is about playing with Nova's bright display modes and learning how a few color choices can turn plain text or simple shapes into something that feels like a real screen from a game.}
\programlisting{Tiny 3}{tiny-3.bas}{draw}{Before a game has sprites, it needs marks on the screen. This listing invites you to draw lines, dots, boxes, and patterns on Nova's larger graphics display, then start asking the fun questions: what happens if the computer draws faster, brighter, or in a different direction?}
\programlisting{Tiny 4}{tiny-4.bas}{sound}{A silent computer feels asleep. This program wakes Nova up with tones, pauses, and little musical ideas, showing how sound can make even a tiny listing feel playful. Type it in, change a number, and you can hear the machine answer back.}
\programlisting{Tiny 5}{tiny-5.bas}{sprite}{This is the first step toward arcade motion. A hardware sprite can move without redrawing the whole screen, which means NovaBASIC can keep the program simple while the video hardware does the busy work. Once one sprite moves, ships, balls, bats, monsters, and players are all within reach.}
\programlisting{Roadhog}{roadhog.bas}{road}{Roadhog is for anyone who wants the screen to feel like it is moving under the player. Nova's display hardware can handle colorful roads, smooth sprite traffic, and timed updates, so this starter points toward a driving game that feels faster and cleaner than old type-in road games ever could.}
\programlisting{Piano}{piano.bas}{piano}{This is a program you can play even before it becomes a game. The idea is simple: turn the keyboard into an instrument, let Nova's sound hardware carry the tune, and then experiment with rhythm, pitch, and recording. It is a friendly doorway into making the machine sing.}
\programlisting{Bets}{bets.bas}{money}{Bets turns numbers into suspense. It is a good place to learn how a game can feel exciting with text, money, odds, and quick feedback instead of lots of graphics. Nova makes the table clean and fast, so the player can focus on the next risky choice.}
\programlisting{Safe}{safe.bas}{safe}{A safe-cracking puzzle is small, but it has everything a good program needs: mystery, memory, feedback, and a reason to try again. Nova can add color and sound to every guess, turning a few numbers into a tense little machine on the screen.}
\programlisting{Boswain}{boswain.bas}{ship}{Boswain aims for action on a busy deck. The fun is in steering through trouble while the screen shows clear movement and the sound gives each mistake some bite. Nova's sprites make it a natural fit for a fast nautical challenge instead of a slow text-only exercise.}
\programlisting{Hanoi}{hanoi.bas}{tower}{Tower of Hanoi is a classic because the rules are tiny and the puzzle grows in your head. Nova lets the program draw the disks, move them, and make each step visible, so a brain teaser becomes something you can watch, solve, and improve.}
\programlisting{Miser}{miser.bas}{money}{Miser is about squeezing choices out of limited resources. It is the kind of game where a plain number can start to feel precious, and NovaBASIC is perfect for making those turns readable, fast, and easy to change into your own strategy game.}
\programlisting{Mad}{mad.bas}{words}{Mad is proof that games do not need explosions to be fun. With strings, prompts, and a little surprise, Nova can become a story machine that makes people laugh at what they typed a minute earlier. It is also a gentle way to learn how programs handle words.}
\programlisting{Godzilla}{godzilla.bas}{monster}{This one should feel big, loud, and a little ridiculous. Nova's sprites and graphics can turn a simple city grid into something stompable, with buildings to smash and sound effects to make the damage feel satisfying. It is a great excuse to make the computer misbehave on purpose.}
\programlisting{Ratrun}{ratrun.bas}{maze}{Maze games are where speed and clarity matter. Ratrun points toward quick tile drawing, crisp movement, and a chase that uses Nova's larger screen to show more of the maze at once. Type it in and you are halfway to building your own arcade labyrinth.}
\programlisting{Yahtzee}{yahtzee.bas}{dice}{Dice games are perfect for learning how a program manages turns, choices, and score sheets. Nova can make the dice and scoring grid easy to read, while BASIC keeps the rules close enough that you can tinker with them without getting lost.}
\programlisting{Leap}{leap.bas}{sprite}{Leap is a one-button kind of idea, which makes it dangerous in the best way. The player waits, jumps, misses, tries again, and soon the timing gets personal. Nova's frame timing and sprite motion can make that simple loop feel sharp and fair.}
\programlisting{Box}{box.bas}{board}{Box is for puzzle builders. A board, a few pieces, and one clever rule can keep a player busy for ages. Nova's graphics make the board visible and friendly, while the listing stays small enough that you can redesign the puzzle after you understand it.}
\programlisting{Lawn}{lawn.bas}{board}{Lawn starts with a silly promise: mow the field, do not get clobbered. That gives you a screen that changes as you play, moving hazards, and a clear goal. Nova's sprites and fast redraws can make the grass, mower, and trouble easy to follow.}
\programlisting{Kalah}{kalah.bas}{board}{Kalah is an ancient-feeling board game with modern programming lessons hiding inside it. You get pits, stones, turns, captures, and simple strategy. Nova can draw the board and counters clearly, making it easier to understand both the game and the code.}
\programlisting{Bonzo}{bonzo.bas}{sprite}{Bonzo is a place for character. Once a little figure moves on screen, every added frame, sound, and obstacle makes it feel more like a real arcade game. Nova's sprite hardware lets BASIC spend less time drawing and more time deciding what Bonzo should do next.}
\programlisting{Bjack}{bjack.bas}{dice}{Blackjack makes arithmetic feel like risk. The program deals, the player decides, and every card changes the mood of the hand. Nova can keep the table readable and lively, with enough polish that a text-and-card game still feels like a night at the arcade counter.}
\programlisting{Fire}{fire.bas}{fire}{Fire is about pressure. Something is spreading, time is short, and the player has to act before the screen gets out of control. Nova's color, sprites, and sound can make danger visible right away, which is exactly what an action listing needs.}
\programlisting{Zap}{zap.bas}{target}{Zap is pure reflex fun: see the target, move fast, fire, and hear the result. Hardware sprites are made for this kind of quick action, and Nova lets BASIC build a target game that feels snappy without burying the reader under complicated drawing code.}
\programlisting{Everest}{everest.bas}{mountain}{Everest should feel like a climb, not just a row of numbers. Nova's graphics can show the mountain, the route, the weather, and the danger, while the program turns each decision into another step upward. It is a survival story you can type.}
\programlisting{Reversi}{reversi.bas}{board}{Reversi is a quiet game that becomes fierce once the board starts flipping. This listing direction shows how Nova can make a clean 8 by 8 board, highlight moves, and redraw pieces fast enough that strategy stays in the foreground.}
\programlisting{Bop}{bop.bas}{sound}{Bop is all about rhythm and reaction. A beep, a flash, a pause, and suddenly the player is leaning toward the keyboard waiting for the next cue. Nova's sound and timing make it a perfect little laboratory for building games that feel musical.}
\programlisting{Spot}{spot.bas}{target}{Spot turns the screen into a memory and observation test. Colored cells, quick changes, and simple scoring can create a surprising amount of tension. Nova's bright display makes patterns easy to see, which means the challenge can come from the game instead of the screen being muddy.}
\programlisting{Dots}{dots.bas}{board}{Dots starts as a school-paper game and becomes a programming lesson in lines, boxes, and turns. Nova can draw the grid cleanly, update only what changed, and make each captured square feel earned. It is a small game with a very visible payoff.}
\programlisting{Capture}{capture.bas}{board}{Capture is for players who like taking over space. A board game about territory gives NovaBASIC a chance to show fast redraws, clear ownership, and computer-managed rules. Once the board is on screen, it begs for new strategies and house rules.}
\programlisting{Dive}{dive.bas}{water}{Dive should feel like motion through water: timing the drop, judging the path, and chasing a better score. Nova's sprites can handle the diver while graphics make the pool or sea feel alive, giving a small arcade idea room to splash around.}
\programlisting{Stop}{stop.bas}{target}{Stop is the kind of program that looks easy until you try to beat it. A marker moves, your timing matters, and the machine judges you instantly. Nova's frame-based timing keeps the challenge honest, so improving feels like skill rather than luck.}
\programlisting{Piegram}{piegram.bas}{chart}{Piegram turns math into a picture. It is not just a chart program; it is a chance to see angles, slices, colors, and labels become something visual. Nova's graphics make it a good first step toward dashboards, gauges, and game status screens.}
\programlisting{Bat}{bat.bas}{paddle}{Bat is the doorway to paddle games. A moving ball, a player-controlled bat, and simple collision checks can become tennis, breakout, or a dozen other arcade ideas. Nova's sprites and timing do the heavy lifting so the listing can stay readable.}
\programlisting{Rescue}{rescue.bas}{rescue}{Rescue gives the player a job that matters: get in, save someone, and get out before things go wrong. Nova's sprites, scrolling pressure, and sound can make a small mission feel urgent, which is exactly the kind of magic a type-in game should deliver.}
